\chapter{发电厂的生产系统及热经济性}

\begin{problem}
凝汽式火电厂热力系统主要包含哪些局部热力系统？各局部热力系统的设备组成及主要作用是什么？
\end{problem}

\begin{solution}
发电厂全面性热力系统一般由主蒸汽和再热蒸汽系统、旁路系统、回热加热（回热抽汽及疏水）系统、给水系统、除氧系统、主凝结水系统、补充水系统、锅炉排污系统、供热系统、厂内循环水系统和锅炉启动系统等局部系统组成。PPT中详细介绍了以下局部热力系统的组成及作用：

\begin{enumerate}
  \item \textbf{主蒸汽及再热蒸汽系统}：一次蒸汽（主蒸汽）系统由从锅炉过热器出口至汽轮机进口的主蒸汽管道、阀门及疏水装置和通往各用新蒸汽的支管组成；二次蒸汽（再热蒸汽）系统由高压缸排汽至锅炉再热器入口的冷再热管道、阀门，和从再热器出口至中压缸进口的热再热管道、阀门组成。
  \item \textbf{旁路系统}：由与汽轮机并联的减压减温器及相连管道组成。其作用包括保护再热器；协调启动参数和流量，缩短启动时间，延长汽轮机寿命；回收工质和热量、降低噪声；防止锅炉超压；在电网故障或甩负荷时维持锅炉热备用。
  \item \textbf{给水系统}：设备通常包含吸水母管、压力母管和锅炉给水母管及其切换阀门等。其作用是输送流量大、压力高的工质，对发电厂的安全、经济、灵活运行至关重要。
\end{enumerate}
\end{solution}

\begin{problem}
何谓发电厂的原则性和全面性热力系统？二者的区别何在？图9-5所示原则性热力系统所反映的设备构成连接有何特点？
\end{problem}

\begin{solution}
\begin{enumerate}
  \item \textbf{原则性热力系统}：只表示热力设备之间的本质联系，表明能量转换与利用的基本过程，反映动力循环中工质的基本流程和完善程度。
  \item \textbf{全面性热力系统}：是在原则性热力系统的基础上，充分考虑到发电厂生产所必须的连续性、安全性、可靠性和灵活性后所组成的实际热力系统。
  \item \textbf{区别}：全面性热力系统如实反映所有主辅设备（包括备用设备）的数量和联系，所有管道附件全部表示于图中；而原则性热力系统相同设备只表示一个，备用设备不予表示，设备之间的联系以单线表示，管道附件一般也不表示。这是二者在画法上的根本区别。
  \item \textbf{原则性热力系统的设备构成连接特点}：（结合原则性热力系统的定义）其特点是高度抽象化，同类设备仅画出一个代表，不显示备用系统与繁杂附件，管线连接全部采用单线，直观突出了热力循环的本质流程与能量转换路径。
\end{enumerate}
\end{solution}

\begin{problem}
火力发电厂的循环供水系统和直流供水系统的区别何在？一般各在什么地区采用？
\end{problem}

\begin{solution}
\begin{enumerate}
  \item \textbf{区别}：直流供水系统是从水源直接取水，吸收乏汽热量后直接返回到原自然水源中；循环供水系统则是让吸热后的冷却水进入冷却塔、喷水池等设施冷却放热，温度降低后再送入凝汽器重复循环使用。
  \item \textbf{适用地区}：直流供水系统要求电厂附近应有充足的水源（流量一般应超过电厂用水量的2-3倍），一般在靠近江河、海洋的地区采用；循环供水系统一般适用于水源不足的内陆地区，其中冷却塔系统尤其适用于气温高或场地受到限制的场合。
\end{enumerate}
\end{solution}

\begin{problem}
凝汽式火力发电厂从燃料化学能转变为电能的整个过程中，可以分为哪些环节？各环节分别在什么设备中完成的？评价各环节能量传递或转换过程完善程度的热经济指标分别是什么？
\end{problem}

\begin{solution}
整个能量转换过程分为以下五个环节：
\begin{enumerate}
  \item \textbf{锅炉设备环节}：在锅炉中完成，将燃料化学能转换为蒸汽热能，评价指标为\textbf{锅炉效率（\(\eta_b\)）}。
  \item \textbf{管道传递环节}：在主蒸汽及再热蒸汽管道中完成，将蒸汽热能传递至汽轮机，评价指标为\textbf{管道效率（\(\eta_p\)）}。
  \item \textbf{汽轮机设备环节}：在汽轮机内完成，将蒸汽热能转换为机械能，评价指标为\textbf{汽轮机内效率（\(\eta_i\)）}。
  \item \textbf{汽轮机输出环节}：在汽轮机主轴输出端完成，进行机械能的传递，评价指标为\textbf{机械效率（\(\eta_m\)）}。
  \item \textbf{发电机环节}：在发电机中完成，将机械能转换为最终的电能，评价指标为\textbf{发电机效率（\(\eta_g\)）}。
\end{enumerate}
\end{solution}

\begin{problem}
凝汽式发电厂的热经济指标包括哪些？各指标所表明的意义是什么？
\end{problem}

\begin{solution}
凝汽式发电厂的热经济指标分为汽轮发电机组指标、锅炉设备指标、主蒸汽及再热蒸汽管道指标和全厂热经济指标四部分。全厂性核心指标及其意义如下：
\begin{enumerate}
  \item \textbf{全厂汽耗量和汽耗率}：表示发电厂生产 \qty{1}{kW \cdot h} 电能所消耗的蒸汽量。
  \item \textbf{全厂热耗量和热耗率}：表示发电厂生产 \qty{1}{kW \cdot h} 电能所消耗的燃料热量。
  \item \textbf{全厂煤耗率和标准煤耗率}：表示生产 \qty{1}{kW \cdot h} 电能所消耗的燃料（或折算的标煤）物理量，常用于不同电厂间的直观比较。
  \item \textbf{全厂热效率}：反映了发电厂输入燃料总热量被转化为有效电能的综合利用程度。
  \item \textbf{发电设备年利用小时数}：衡量发电设备是否充分发挥其应有作用，反映机组在设计、运行维护管理等方面的水平。
\end{enumerate}
\end{solution}

\begin{problem}
发电效率与供电效率的区别何在？试列出供电效率与发电效率及厂用电率之间的关系式。
\end{problem}

\begin{solution}
\begin{enumerate}
  \item \textbf{区别}：发电效率（即全厂热效率）衡量的是燃料总热量转化为发电机发出的\textbf{总电能}的比例；而供电效率则扣除了发电厂生产过程中自身消耗的电能（厂用电），衡量的是燃料总热量转化为实际\textbf{向外供出电能}的比例。
  \item \textbf{关系式}：设全厂热效率（发电效率）为 \(\eta_{pl}\)，供电效率为 \(\eta_{pl}^n\)，厂用电率为 \(K\)，则关系式为：
  \begin{equation*}
    \eta_{pl}^n = \eta_{pl}(1 - K)
  \end{equation*}
\end{enumerate}
\end{solution}

\begin{problem}
供电标准煤耗率和发电标准煤耗率的区别是什么？全厂热效率的高低对煤耗率的影响是怎样的？
\end{problem}

\begin{solution}
\begin{enumerate}
  \item \textbf{区别}：煤耗率是指生产单位电能消耗的燃料量。发电标准煤耗率以发电机发出的总电能为基准进行折算；而供电标准煤耗率则以扣除厂用电后的净供电量为基准进行折算，它更能真实地反映电厂向外界输送可用电能时的燃料消耗水平。
  \item \textbf{对煤耗率的影响}：根据公式 \(b^s = 123 / \eta_{pl}\) [g标准煤/(kW·h)]，标准煤耗率（\(b^s\)）与全厂热效率（\(\eta_{pl}\)）成反比关系。全厂热效率越高，能量利用越完善，生产同样电能所需的煤耗率就越低。
\end{enumerate}
\end{solution}
